
\textbf{Exemples 2.2.}~: a) On peut donner toutes les structures de pré-@@M0@@-anneau de @@M1@@. Il suffit pour cela de se donner @@M2@@, qui sera donc une série formelle à coefficients entiers, soumise à la seule condition de commencer par @@M3@@.

b) Soit @@M4@@ un topos commutativement annelé, et soit @@M5@@ le « groupe de Grothendieck » des @@M6@@-Modules localement isomorphes à un facteur direct d'un @@M7@@-Module libre de type fini~; on sait que @@M8@@ peut être muni d'une structure d'anneau grâce au produit tensoriel (Exp. IV), et nous allons le munir d'une structure de pré-@@M9@@-anneau de la façon suivante~: pour tout @@M10@@-Module @@M11@@, localement facteur direct d'un Module libre de type fini, on désigne par @@M12@@ la classe dans @@M13@@ de la @@M14@@-ième puissance extérieure de @@M15@@ (qui est également localement facteur direct d'un Module libre de type fini)~; ceci définit une fonction sur la catégorie des @@M16@@-Modules localement facteurs directs d'un libre de type fini, à valeurs dans le groupe abélien @@M17@@~; il faut alors montrer que cette fonction est additive, ce qui entraînera qu'elle se factorise par le @@M18@@, d'où la structure cherchée.

\textbf{Lemme 2.2.1.} Soit @@M19@@ \uline{une suite exacte de} @@M20@@-\uline{Modules qui est localement scindée} (\uline{par exemple telle que} @@M21@@ \uline{soit localement facteur direct d'un Module libre de type fini}). \uline{Il existe sur} @@M22@@ \uline{une filtration finie}, \uline{et des isomorphismes}~:
@@M23@@

On prend pour @@M24@@ l'image du morphisme canonique~:
@@M25@@
